\documentclass[a0,portrait]{a0poster}

\usepackage[brazil]{babel}
\usepackage[utf8]{inputenc}
\usepackage[T1]{fontenc}
\usepackage{graphicx}
\usepackage{float}
\usepackage{caption}
\usepackage{xcolor}
\usepackage{tcolorbox}
\usepackage{enumitem}
\usepackage[absolute]{textpos}

% ---- Cores dos graficos do projeto ----
\definecolor{plotblue}{HTML}{3D79F3}
\definecolor{plotred}{HTML}{E6352F}
\definecolor{plotyellow}{HTML}{F9B90A}
\definecolor{plotgreen}{HTML}{34A74B}
\definecolor{plotgray}{HTML}{909090}
\definecolor{flamengored}{HTML}{C8102E}
\definecolor{fluminensegreen}{HTML}{00613C}
\definecolor{fluminensegrenat}{HTML}{7A263A}

% ---- Tipografia ----
\renewcommand{\baselinestretch}{1.3}
\let\normalfontBK\normalfont
\renewcommand\normalfont{\normalfontBK\sffamily}
\setlength{\parskip}{0.35\baselineskip}

\pagestyle{empty}
\setcounter{secnumdepth}{0}

\TPGrid[40mm,40mm]{13}{18}
\parindent=0pt
\parskip=0.4\baselineskip

\let\Textsize\normalsize
\def\LHead#1{\noindent{\LARGE\bf\color{plotgreen} #1}\smallskip}
\def\Title#1{\noindent{\veryHuge #1}}

\newtcolorbox{estat}[1]{
  colback=#1!35, colframe=black,
  arc=10pt, boxrule=1.5pt,
  left=4pt, right=4pt, top=6pt, bottom=6pt
}

\newtcolorbox{destaque}[1]{
  colback=#1!15, colframe=#1,
  arc=8pt, boxrule=3pt,
  left=10pt, right=10pt, top=5pt, bottom=5pt
}

\newcommand{\timefaltas}[5]{
\begin{destaque}{#2}
{\Large\bfseries #1}\\[1mm]
{\large #3}\\[0.5mm]
{\large #4}\\[0.5mm]
{\large #5}
\end{destaque}
}

\newcommand{\timefaltascinza}[4]{
\begin{tcolorbox}[
  colback=plotgray!12, colframe=black,
  arc=8pt, boxrule=3pt,
  left=10pt, right=10pt, top=5pt, bottom=5pt
]
{\Large\bfseries #1}\\[1mm]
{\large #2}\\[0.5mm]
{\large #3}\\[0.5mm]
{\large #4}
\end{tcolorbox}
}

\newcommand{\timefaltasbicolor}[6]{
\begin{tcolorbox}[
  colback=#2!15, colframe=#3,
  arc=8pt, boxrule=3pt,
  left=10pt, right=10pt, top=5pt, bottom=5pt
]
{\Large\bfseries #1}\\[1mm]
{\large #4}\\[0.5mm]
{\large #5}\\[0.5mm]
{\large #6}
\end{tcolorbox}
}

\newcommand{\faltaitem}[2]{%
\begin{minipage}[t]{\linewidth}
\raggedright
\hangindent=10mm
\hangafter=1
{\large\bfseries\color{plotgreen}\makebox[8mm][r]{#1.}}\hspace{1.5mm}{\large #2}
\end{minipage}%
}

\newcommand{\arbitrofaltas}[4]{%
\begin{minipage}[t]{\linewidth}
\raggedright
{\large\bfseries #1}\\[0.5mm]
{\normalsize #2}\\[0.5mm]
{\normalsize #3}\\[0.5mm]
{\normalsize #4}
\end{minipage}%
}

\begin{document}
\sffamily
\baselineskip=1.15\baselineskip


%%----------------------------------------------------------
%% TITULO
%%----------------------------------------------------------
\begin{textblock}{13}(0, 0.3)
\begin{center}
\Title{Pr\'atica em Ci\^encia de Dados III}\\
\vspace{2mm}
{\huge\bfseries\color{plotgreen}
Classificando motivos de Cart\~oes do Futebol Brasileiro}\\
\vspace{4mm}
{\LARGE Gabriel Carbinatto \hspace{5mm}|\hspace{5mm}
Arthur Martins Ferreira de Sousa}\\
\vspace{1mm}
{\Large Universidade de S\~ao Paulo ---
Instituto de Ci\^encias Matem\'aticas e de Computa\c{c}\~ao}
\end{center}
\end{textblock}


%%----------------------------------------------------------
%% NUMEROS EM DESTAQUE
%%----------------------------------------------------------
\begin{textblock}{2.8}(0.2, 2.0)
\begin{estat}{plotblue}
\begin{center}
{\veryHuge\bfseries\color{black} 5.681}\\[1mm]
{\Large\color{black} jogos analisados}
\end{center}
\end{estat}
\end{textblock}

\begin{textblock}{2.8}(3.5, 2.0)
\begin{estat}{plotred}
\begin{center}
{\veryHuge\bfseries\color{black} 25.838}\\[1mm]
{\Large\color{black} cart\~oes estudados}
\end{center}
\end{estat}
\end{textblock}

\begin{textblock}{2.8}(6.8, 2.0)
\begin{estat}{plotgreen}
\begin{center}
{\veryHuge\bfseries\color{black} 18}\\[1mm]
{\Large\color{black} tipos de falta}
\end{center}
\end{estat}
\end{textblock}

\begin{textblock}{2.8}(10.1, 2.0)
\begin{estat}{plotyellow}
\begin{center}
{\veryHuge\bfseries\color{black} 7 anos}\\[1mm]
{\Large\color{black} 2020--2026}
\end{center}
\end{estat}
\end{textblock}


%%----------------------------------------------------------
%% COL ESQUERDA 1 — Os Dados + O Problema
%%----------------------------------------------------------
\begin{textblock}{6.0}(0, 3)

\LHead{Os Dados Coletados}

Coletamos \textbf{5.681 s\'umulas da CBF} de \textbf{3 competi\c{c}\~oes},
totalizando \textbf{25.838 cart\~oes} analisados ao longo de \textbf{7 anos}:

\vspace{2mm}
\begin{destaque}{plotblue}
\Large
\begin{tabular}{l l}
  \textbf{S\'erie A} (1$^a$ divis\~ao) & 2020--2026 \\[2mm]
  \textbf{S\'erie B} (2$^a$ divis\~ao) & 2020--2026 \\[2mm]
  \textbf{Copa do Brasil}               & 2020--2026 \\
\end{tabular}
\end{destaque}

\vspace{5mm}
\LHead{O Problema}

Toda partida de futebol tem um documento oficial chamado \textbf{s\'umula}.
Nele, o \'arbitro anota o motivo de cada cart\~ao com suas \textbf{pr\'oprias palavras}!

Veja como a \emph{mesma falta} pode aparecer de formas diferentes:

\vspace{2mm}
\begin{destaque}{plotyellow}
{\large\bfseries\color{plotgray} A\c{c}\~ao temer\'aria:}\\[1mm]
\itshape\large
``Por cal\c{c}ar seu advers\'ario de maneira temer\'aria'' \\[1.5mm]
``Entrada temer\'aria sobre o advers\'ario'' \\[1.5mm]
``Por cometer entrada temer\'aria''
\end{destaque}

\vspace{2mm}
\begin{destaque}{plotred}
{\large\bfseries\color{plotgray} Protesto contra a arbitragem:}\\[1mm]
\itshape\large
``Por reclamar ostensivamente com palavras e gestos.'' \\[1.5mm]
``Desaprovar com palavras e gestos a decis\~ao da arbitragem.'' \\[1.5mm]
``Ap\'os reclamar de forma acintosa das decis\~oes da arbitragem.''
\end{destaque}

\vspace{2mm}
\begin{destaque}{plotgreen}
{\large\bfseries\color{plotgray} Falta t\'atica em ataque promissor:}\\[1mm]
\itshape\large
``Por impedir um ataque promissor.'' \\[1.5mm]
``Por cal\c{c}ar seu advers\'ario impedindo um ataque promissor.'' \\[1.5mm]
``Segurar o advers\'ario impedindo um ataque promissor.''
\end{destaque}

\vspace{2mm}
\begin{destaque}{plotblue}
{\large\bfseries\color{plotgray} Toque de m\~ao na bola:}\\[1mm]
\itshape\large
``Por tocar a bola com a m\~ao de forma deliberada.'' \\[1.5mm]
``Por interceptar a bola com o bra\c{c}o.'' \\[1.5mm]
``Por colocar a m\~ao na bola impedindo a jogada.''
\end{destaque}

\end{textblock}


%%----------------------------------------------------------
%% COL ESQUERDA 2 — Os 18 Tipos de Falta
%%----------------------------------------------------------
\begin{textblock}{6.0}(0, 12.5)

\LHead{Os 18 Tipos de Falta Descobertos}

\vspace{2mm}
\renewcommand{\arraystretch}{1.35}
\begin{tabular}{@{}p{0.48\linewidth}p{0.48\linewidth}@{}}
\faltaitem{1}{A\c{c}\~ao temer\'aria} &
\faltaitem{10}{Conduta violenta} \\
\faltaitem{2}{Protesto contra arbitragem} &
\faltaitem{11}{Jogo brusco grave} \\
\faltaitem{3}{Falta t\'atica} &
\faltaitem{12}{Toque de m\~ao na bola} \\
\faltaitem{4}{Retardar rein\'icio} &
\faltaitem{13}{Impedir clara oportunidade de gol} \\
\faltaitem{5}{Agarrar advers\'ario} &
\faltaitem{14}{Insubordina\c{c}\~ao t\'ecnica} \\
\faltaitem{6}{Conduta antidesportiva} &
\faltaitem{15}{Simula\c{c}\~ao} \\
\faltaitem{7}{Confronto f\'isico} &
\faltaitem{16}{N\~ao respeitar dist\^ancia} \\
\faltaitem{8}{Comemora\c{c}\~ao irregular} &
\faltaitem{17}{Infringir regras repetidas} \\
\faltaitem{9}{Gesto/linguagem ofensiva} &
\faltaitem{18}{Entrar/sair sem permiss\~ao} \\
\end{tabular}

\end{textblock}


%%----------------------------------------------------------
%% COL DIREITA 1 — Como foi feito
%%----------------------------------------------------------
\begin{textblock}{6.0}(7, 3)

\LHead{Como foi feito?}

\begin{destaque}{plotblue}
{\large\bfseries Passo 1 --- Extra\c{c}\~ao}\\[1mm]
\large Baixamos as s\'umulas da CBF e organizamos os dados dos jogos
\end{destaque}

\vspace{2mm}
\begin{destaque}{plotblue}
{\large\bfseries Passo 2 --- Leitura}\\[1mm]
\large Extra\'imos dos PDFs o motivo escrito para cada cart\~ao
\end{destaque}

\vspace{2mm}
\begin{destaque}{plotblue}
{\large\bfseries Passo 3 --- Entender}\\[1mm]
\large Transformamos cada justificativa em uma representa\c{c}\~ao sem\^antica
\end{destaque}

\vspace{2mm}
\begin{destaque}{plotblue}
{\large\bfseries Passo 4 --- Agrupar}\\[1mm]
\large Textos com significado parecido foram reunidos automaticamente
\end{destaque}

\vspace{2mm}
\begin{destaque}{plotblue}
{\large\bfseries Passo 5 --- Validar}\\[1mm]
\large Verificamos se grupos encontrados estavam corretamente agrupados
\end{destaque}

\end{textblock}


%%----------------------------------------------------------
%% COL DIREITA 2 — Destaques por time
%%----------------------------------------------------------
\begin{textblock}{6.0}(7, 7.5)

\LHead{Quais motivos se destacam nos times?}

\vspace{2mm}
\begin{minipage}[t]{0.48\linewidth}
\timefaltas{Palmeiras}{plotgreen}
{Gesto/linguagem ofensiva}
{Comemora\c{c}\~ao irregular}
{Protesto contra arbitragem}
\end{minipage}
\hfill
\begin{minipage}[t]{0.48\linewidth}
\timefaltascinza{Corinthians}
{Comemora\c{c}\~ao irregular}
{Confronto f\'isico}
{A\c{c}\~ao temer\'aria}
\end{minipage}

\vspace{2mm}
\begin{minipage}[t]{0.48\linewidth}
\timefaltas{S\~ao Paulo}{plotred}
{Confronto f\'isico}
{Protesto contra arbitragem}
{Insubordina\c{c}\~ao t\'ecnica}
\end{minipage}
\hfill
\begin{minipage}[t]{0.48\linewidth}
\timefaltascinza{Santos}
{Comemora\c{c}\~ao irregular}
{Agarrar advers\'ario}
{Confronto f\'isico}
\end{minipage}

\vspace{2mm}
\begin{minipage}[t]{0.48\linewidth}
\timefaltasbicolor{Flamengo}{flamengored}{black}
{A\c{c}\~ao temer\'aria}
{Comemora\c{c}\~ao irregular}
{Confronto f\'isico}
\end{minipage}
\hfill
\begin{minipage}[t]{0.48\linewidth}
\timefaltasbicolor{Fluminense}{fluminensegreen}{fluminensegrenat}
{Confronto f\'isico}
{Protesto contra arbitragem}
{Conduta antidesportiva}
\end{minipage}

\vspace{5mm}
\LHead{\'Arbitros: motivos que se destacam}

\vspace{1.5mm}
\renewcommand{\arraystretch}{1.25}
\begin{tabular}{@{}p{0.32\linewidth}p{0.32\linewidth}p{0.32\linewidth}@{}}
\arbitrofaltas{Anderson Daronco}{Falta t\'atica}{Confronto f\'isico}{Conduta violenta} &
\arbitrofaltas{Ramon Abatti}{Agarrar advers\'ario}{Comemora\c{c}\~ao irregular}{Confronto f\'isico} &
\arbitrofaltas{Raphael Claus}{Conduta antidesportiva}{A\c{c}\~ao temer\'aria}{Jogo brusco grave}
\end{tabular}

\vspace{4mm}

\begin{tabular}{@{}p{0.32\linewidth}p{0.32\linewidth}p{0.32\linewidth}@{}}
\arbitrofaltas{Wilton Pereira}{Gesto/linguagem ofensiva}{Protesto contra arbitragem}{Conduta antidesportiva} &
\arbitrofaltas{Braulio da Silva}{Gesto/linguagem ofensiva}{Toque de m\~ao na bola}{Insubordina\c{c}\~ao t\'ecnica} &
\arbitrofaltas{Bruno Arleu}{A\c{c}\~ao temer\'aria}{Confronto f\'isico}{Jogo brusco grave}
\end{tabular}

\end{textblock}


%%----------------------------------------------------------
%% COL DIREITA 3 — Conclusao
%%----------------------------------------------------------
\begin{textblock}{6.0}(7, 14)

\LHead{Conclus\~ao}

\Large
A ci\^encia de dados ajudou a transformar milhares de s\'umulas em
algo muito mais f\'acil de entender.

Em vez de ler cart\~ao por cart\~ao, usamos o computador para encontrar
frases parecidas, organizar os motivos em grupos e mostrar padr\~oes que
aparecem nos jogos.

Assim, um problema grande e repetitivo virou uma vis\~ao clara sobre
o que acontece dentro de campo.

\end{textblock}


%%----------------------------------------------------------
%% BARRA + LOGO CENTRALIZADA
%%----------------------------------------------------------
\begin{textblock}{13.5}(-0.3, 17.1)
\rule{\textwidth}{2pt}
\end{textblock}

\begin{textblock}{13}(0, 17.3)
\begin{center}
\includegraphics[width=14cm]{viz/logo_ICMC (4).png}
\end{center}
\end{textblock}


\end{document}
